\subsection{What the opened literature establishes}

\begin{statement}{PR.1}{Known}{The moduli-group interpretation supplies an inclusion}
With inertia preservation and geometric symmetric-group compatibility included,
\citet[Main Theorem and Corollary]{LitHarbaterSchneps2000} identify
the four-point outer automorphism group with $\GZ$ and the corresponding
$n$-point group, $n\ge5$, with $\GT$. The forgetful map supplies
$\GT\hookrightarrow\GZ$.
\end{statement}
\begin{proof}[Derivation and scope]
The four-point pure fundamental group is the marked $\F$ of the prompt.
Its three peripheral classes are $x,y,(xy)^{-1}$, and its geometric
symmetric action is the specified $S_3$. The source's Proposition 4
turns the cyclic compatibility into the three-cycle equation; the
transposition compatibility gives the two-cycle equation. At five
points the remaining compatibility is the pentagon. Restricting an
actual compatible automorphism along the forgetting maps yields the
injection in the source's corollary. This identifies the groups and
maps used here, but the cited theorem does not assert that the injection
is proper. The proof of properness in P.10 below uses the supplied R1
and independently constructed pro-prime inputs.
\end{proof}

\begin{statement}{PR.2}{Known}{Categorical coherence and its faithful scope}
Drinfeld's group pentagon is the universal associativity coherence
condition for the modification by a pair $(\lambda,f)$; a single
pointed braided category need not detect that group relation.
\end{statement}
\begin{proof}[Derivation]
In \citet[Section 4, equations (4.5), (4.10), (4.12), (4.13)]{LitDrinfeld1991},
the associator changes by evaluation of $f$ on the two pure monodromies,
and the braiding changes by the multiplier. The free braid situation
retains the group equation. Passing to a category applies its braid
representation to the equation, which is a valid implication in that
direction. The representation may have a kernel. In the category used
in A.2(c), the two relevant monodromies commute, so every $f\in\F'$
evaluates to one; the explicit proof and finite coherence checks below
show that this loss of information occurs here. Thus universal
coherence cannot be inferred back from this one example.
\end{proof}

\begin{statement}{PR.3}{Known}{The classical torsion result used here is a centralizer theorem}
For $c=(-1,1)$, \citet[Proposition 4(ii)]{LitLochakSchneps1997}
prove $C_{\GT}(c)=\langle c\rangle$. This result does not itself
classify every torsion element of $\GT$ or $\GZ$. We use the supplied,
refereed R7 for the assertion that every nonidentity finite-order
element of either group is an involution of multiplier $-1$.
\end{statement}
\begin{proof}[Derivation and separation of assertions]
The source's $c$ is the pair with inverted peripheral generators and
trivial $f$, matching our marked normalization. Its centralizer assertion
therefore transfers literally. A centralizer statement at this element
cannot imply that a different torsion element is conjugate to it.
R7 separately gives torsion-freeness of $\HH$, the order-two conclusion,
and $\Delta((-1,f))=\delta((1,f))$. R4 gives zero on every conjugate
of $c$, because $\Delta$ is a character and $c$ is Galois. No general
involution-conjugacy theorem is claimed or needed in this return.
\end{proof}

\begin{statement}{PR.4}{Known}{The currently available mapping-class theorems retain comparison hypotheses}
The braid, configuration-space, and higher-genus automorphism theorems
listed below do not supply the CM comparison needed to evaluate
$\Delta|_{\GT}$.
\end{statement}
\begin{proof}[Exact reductions]
\citet[Theorems A and B]{LitMinamideNakamura2022} treat profinite braid
groups on at least four strands and their central quotients. Their
additional central-unit factor disappears only after that specified
quotient; the theorem cannot be specialized to the three-strand modular
source. \citet[Corollaries B and C]{LitHoshiMinamideMochizuki2022}
recover the configuration-space fibre structure and the GT factor in
the outer automorphism group under their stated number-of-points and
prime-set hypotheses. Identification of that factor gives no equality
with Galois or calculation of its abelianization.

\citet[Theorem 1.1]{LitLochakNakamuraSchneps2026} construct full
profinite actions in genus $g>0$ with zero or one marked point.
Their Section 8 retains the required cutting and Lego compatibilities
as further work. Therefore the existence of these individual actions
cannot identify two actions defined by different CM constructions.
The exact low-genus sewing application of
\citet[Theorem 2.3 and Proposition 2.6(i)]{LitBoggi2026v6} is proved
in Part B; it is an identification of the relevant groups with their
maps, and supplies no equality of arbitrarily specified actions.
\citet[Theorem A and Corollary B]{BoggiSymmetric2026} give the
congruence property for finite-subgroup centralizers when the quotient
surface has genus at most two. That result controls the induced
profinite topology, not the residual frame character.
\end{proof}

\begin{statement}{PR.5}{Known}{Special-locus results and the exact Galois-only comparison}
The harmonic, equianharmonic and periodic-frame constructions occur
in the opened literature, but none of the specific cited results
identifies the two CM coordinates for every element of $\GT$.
\end{statement}
\begin{proof}[Derivation and quantifiers]
\citet[Theorems 1.1--1.2]{LitNakamuraSchneps2000} define the subgroup
with the extra relations $(\mathrm{III}')$ and $(\mathrm{IV})$ and
prove the Galois image lies in it. The harmonic and equianharmonic
identities are printed in the primary report
\citet[Theorems 1--2]{LitTsunogaiNakamura2002}; the later 2003
article's mathematical full text was not accessible here and is marked
unverified in the bibliography. \citet[Theorem 3.4]{LitSchneps2006}
deduces the further identities for its precisely defined special-locus
subgroup, including $\lambda=1$, $\rho_2=\rho_3=0$ and two additional
compatibilities. Those restrictions cannot be removed by citing the
theorem.

\citet[Proposition 4.4]{LitNakamuraTsunogaiYasuda2010} give an
all-GT triangle action, while Propositions 4.11 and 4.13 and
Theorem 4.16 compare it with the standard action on Galois. The
Galois restriction is explicit; Section 4.12 is not a proposition.
Part B derives its sign with the central powers retained.
\citet[Theorems A and B]{Collas2012GenusZero} give geometricity and
cyclotomic conjugacy for prime-order torsion in genus zero. The
separate genus-one text \citep{CollasGenusOneAuthorText} originally
uses the smaller acting group $I\Gamma$. Neither paper computes the
relative CM parity character of this prompt.
\end{proof}

\begin{statement}{PR.6}{Open}{Search boundary for finite characters, strictness and novelty}
The opened literature search did not locate an independent theorem
proving the precise R1 lifting statement, the R2 prescribed mixed-frame
lift, the R4 relative quadratic character, or the complete R7 torsion
classification. It also did not locate a later solution of the profinite
abelianization question in \citet[Question 1.2]{LitLochakSchnepsOpenProblems}.
This is a record of the search, not a proof that no such source exists.
\end{statement}
We searched both available web indexes for the named sources, exact
titles, profinite versus pro-prime strictness, finite characters,
non-Galois gentle elements, torsion and CM comparisons. The primary
sources and exact loci are recorded in
\texttt{agent\_literature/source\_audit.md}. We checked the revision
of \citet{LitCollas2026v2} dated 5 November 2026; its strong
indecomposability assertion concerns the group rather than its
abelianization. In \citet[Proposition 5.20 and closing remarks,
p.373]{LitGuillot2014}, an unhatted inverse-limit group denotes the
usual $\GZ$; the hatted pentagon group is described as a subgroup,
without an assertion of properness. That notation cannot establish
strictness. R2 and R4 do give an actual non-Galois element of $\GZ$
under the accepted supplied results. We claim no historical priority
for that deduction or for any proof below.
